\section*{Capítulo \ref{ch:b2:genfun} --- Funciones generatrices de probabilidad}

\begin{solution}{exo:b2:genfun:1}
Dado equilibrado: $G(t) = \frac16(t + t^2 + \dots + t^6) = \frac
t6(1 + t + \dots + t^5)$. Si la suma de dos dados equilibrados fuera
uniforme sobre $\{2, \dots, 12\}$, entonces
\[
G(t)^2 = \frac{t^2}{36}\,h(t)^2
= \frac{t^2}{11}\sum_{k=0}^{10}t^k ,
\qquad h(t) = 1 + t + \dots + t^5 .
\]
Ahora bien, $h$ es un polinomio real de grado impar $5$, luego tiene
una raíz real (teorema del valor intermedio; en concreto, $h(-1) =
0$), y por tanto $h^2$ también. Pero $\sum_{k=0}^{10}t^k$ no tiene
ninguna: es positivo para $t \geq 0$ y, para $t < 0$, vale
$\frac{t^{11} - 1}{t - 1}$, un cociente de dos números negativos.
Contradicción; la suma de dos dados equilibrados nunca es uniforme
(como confirma la familiar \omterm{def:b2:randomvar:law}{distribución} triangular de las sumas de
dados).
\end{solution}

\begin{solution}{exo:b2:genfun:2}
\emph{Binomial:} $G(t) = (1 - p + pt)^n$, $G'(t) = np(1 - p +
pt)^{n-1}$, $G''(t) = n(n-1)p^2(1 - p + pt)^{n-2}$, de modo que
\[
\E(X) = G'(1) = np,
\qquad
V(X) = G''(1) + G'(1) - G'(1)^2
= n(n-1)p^2 + np - n^2p^2 = np(1-p).
\]
\emph{Geométrica} ($q = 1 - p$): $G(t) = \frac{pt}{1 - qt}$, luego
$G'(t) = \frac{p}{(1 - qt)^2}$ y $G''(t) = \frac{2pq}{(1 - qt)^3}$;
en $t = 1$ (usando $1 - q = p$):
\[
\E(X) = \frac{p}{p^2} = \frac1p,
\qquad
V(X) = \frac{2q}{p^2} + \frac1p - \frac{1}{p^2}
= \frac{2q + p - 1}{p^2}
= \frac{q}{p^2} ,
\]
que coincide con el~\cref{exo:b2:randomvar:1} con menos trabajo.
\end{solution}

\begin{solution}{exo:b2:genfun:3}
No, ni siquiera con trucajes distintos. Supongamos que $X, Y$ son
\omterm{def:b2:randomvar:law}{leyes} sobre $\{1, \dots, 6\}$ con suma uniforme. Entonces $G_X(t) =
t\,a(t)$ y $G_Y(t) = t\,b(t)$ con $a, b$ polinomios reales de grado
\emph{a lo sumo} $5$; y sus grados han de sumar $10$ (la suma alcanza
$12$ con probabilidad positiva), de modo que $\deg a = \deg b = 5$,
ambos impares. Como en el~\cref{exo:b2:genfun:1},
\[
a(t)\,b(t) = \frac{1}{11}\sum_{k=0}^{10}t^k
\]
forzaría una raíz real en el miembro izquierdo (todo polinomio real
de grado impar tiene una) y ninguna en el derecho. Así pues, ningún
trucaje de dos dados \omterm{def:b2:proba:independence}{independientes} ---iguales o no--- produce una
suma uniforme.
\end{solution}

\begin{solution}{exo:b2:genfun:4}
Por el~\cref{thm:b2:genfun:compound} con $G_N(s) =
e^{\lambda(s-1)}$ y $G_X(t) = 1 - p + pt$:
\[
G_S(t) = e^{\lambda(1 - p + pt - 1)} = e^{\lambda p(t - 1)} :
\]
$S \sim \mathcal{P}(\lambda p)$. El recuento descartado $D = N - S$
cuenta los mismos elementos conservados con probabilidad $1 - p$, así
que, por el mismo cálculo, $D \sim \mathcal{P}(\lambda(1 - p))$.
\omterm{def:b2:proba:independence}{Independencia}, directamente: para $j, k \in \N$,
\begin{align*}
\P(S = j,\ D = k)
&= \P(N = j + k)\,\binom{j+k}{j}p^jq^k
= e^{-\lambda}\frac{\lambda^{j+k}}{(j+k)!}\,
\frac{(j+k)!}{j!\,k!}\,p^jq^k\\
&= \Bigl(e^{-\lambda p}\frac{(\lambda p)^j}{j!}\Bigr)
\Bigl(e^{-\lambda q}\frac{(\lambda q)^k}{k!}\Bigr)
\end{align*}
con $q = 1 - p$: la \omterm{def:b2:randomvar:law}{ley} conjunta factoriza como
$\mathcal{P}(\lambda p) \otimes \mathcal{P}(\lambda q)$. Un flujo de
Poisson dividido al azar rinde flujos de Poisson
\emph{\omterm{def:b2:proba:independence}{independientes}}; un pequeño milagro usado constantemente en teoría
de colas.
\end{solution}

\begin{solution}{exo:b2:genfun:5}
Los tiempos de espera entre caras consecutivas son variables
geométricas $\mathcal{G}(p)$ \omterm{def:b2:proba:independence}{independientes} (ausencia de memoria: tras
cada cara el juego recomienza), de modo que $T_r = W_1 + \dots + W_r$
y la multiplicatividad (\cref{thm:b2:genfun:product}) da
\[
G_{T_r}(t) = \Bigl(\frac{pt}{1 - qt}\Bigr)^{r},
\qquad
\E(T_r) = r\,\E(W_1) = \frac rp,
\qquad
V(T_r) = r\,V(W_1) = \frac{rq}{p^2}
\]
($q = 1 - p$; las \omterm{def:b2:randomvar:variance}{varianzas} se suman por \omterm{def:b2:proba:independence}{independencia}). Desarrollo:
por la serie binomial generalizada (\cref{ch:b2:powerseries}), $(1 -
qt)^{-r} = \sum_{m\geq0} \binom{m + r - 1}{r - 1}q^mt^m$, de modo que
el coeficiente de $t^n$ en $p^rt^r(1 - qt)^{-r}$ es (con $m = n - r$)
\[
\P(T_r = n) = \binom{n-1}{r-1}p^r(1-p)^{n-r},
\qquad n \geq r ,
\]
la \omterm{def:b2:randomvar:law}{ley} \emph{binomial negativa}; combinatoriamente: la $r$-ésima
cara cae en el lanzamiento $n$ si y solo si las $r - 1$ caras
anteriores eligen sus posiciones entre los $n - 1$ primeros
lanzamientos.
\end{solution}

\begin{solution}{exo:b2:genfun:6}
$m = 1\cdot\frac38 + 2\cdot\frac38 + 3\cdot\frac18 = \frac{3 + 6 +
3}{8} = \frac32 > 1$: supercrítico. La \omterm{ex:b2:powerseries:fibonacci}{función generatriz} es
\[
G(t) = \frac{1 + 3t + 3t^2 + t^3}{8} = \frac{(1 + t)^3}{8} ,
\]
así que los puntos fijos resuelven $(1 + t)^3 = 8t$, es decir, $t^3 +
3t^2 - 5t + 1 = 0$. Sacando factor la raíz garantizada $t = 1$:
\[
t^3 + 3t^2 - 5t + 1 = (t - 1)\bigl(t^2 + 4t - 1\bigr),
\]
y $t^2 + 4t - 1 = 0$ da $t = -2 \pm \sqrt5$. La raíz de
$\intco{0}{1}$ es $\sqrt5 - 2 \approx 0.236$: por el~\cref{thm:b2:genfun:extinction},
\[
q = \sqrt 5 - 2 .
\]
(Una comprobación agradable: la \omterm{def:b2:randomvar:law}{ley} de descendencia es la de $3$
monedas equilibradas \omterm{def:b2:proba:independence}{independientes}, $Z_1 \sim \mathcal{B}(3,
\frac12)$.)
\end{solution}

\begin{solution}{exo:b2:genfun:7}
\emph{\omterm{def:b2:randomvar:expectation}{Esperanza}.} Primero, $\E(Z_n) = m^n$: por el~\cref{thm:b2:genfun:compound}, $\E(Z_{n+1}) = \E(Z_n)\,m$, y
$\E(Z_0) = 1$. La familia $\bigl(Z_n(\omega)\P(\{\omega\})
\bigr)_{n, \omega}$ es no negativa, así que Fubini para familias se
aplica incondicionalmente:
\[
\E(Y) = \sum_{n=0}^{\infty}\E(Z_n)
= \sum_{n=0}^\infty m^n = \frac{1}{1 - m} < \infty
\]
(en particular, $Y$ es finita casi seguramente, coherentemente con la
extinción cierta del caso subcrítico).

\emph{Ecuación funcional.} Descompóngase la población según los hijos
del antepasado: si el antepasado tiene $Z_1 = k$ hijos, la
descendencia total es $Y = 1 + Y_1 + \dots + Y_k$, donde $Y_i$ es la
descendencia total del linaje del $i$-ésimo hijo; y las $Y_i$ son
copias independientes de $Y$, independientes de $Z_1$ (linajes
distintos usan \omterm{def:b2:proba:space}{sucesos} de reproducción disjuntos e \omterm{def:b2:proba:independence}{independientes}).
Condicionando a $Z_1$ como en el~\cref{thm:b2:genfun:compound}:
\[
H(t) = \E\bigl(t^Y\bigr)
= t\sum_{k=0}^\infty \P(Z_1 = k)\,H(t)^k
= t\,G\bigl(H(t)\bigr),
\]
donde el factor $t$ da cuenta del propio antepasado. (Para la \omterm{def:b2:randomvar:law}{ley}
$p_0 = 1 - p$, $p_2 = p$ de la ramificación binaria, esta ecuación
cuadrática en $H$ puede resolverse explícitamente y desarrollarse:
los \omterm{ex:b2:powerseries:catalan}{números de Catalan} del~\cref{ch:b2:powerseries} cuentan los
árboles genealógicos.)
\end{solution}

\begin{solution}{exo:b2:genfun:8}
Sea $R > 1$ el radio y fíjese $t \in \intoo{1}{R}$. Entonces
$\E(t^X) = G(t) < \infty$, y la desigualdad de Markov
(\cref{thm:b2:randomvar:markov}) aplicada a la variable no negativa
$t^X$ al nivel $t^n$:
\[
\P(X \geq n) = \P\bigl(t^X \geq t^n\bigr)
\leq \frac{G(t)}{t^n} = C\rho^n,
\qquad C = G(t),\quad \rho = \frac1t \in \intoo{0}{1}.
\]
\emph{Recíproco:} si $\P(X \geq n) \leq C\rho^n$, entonces $p_n \leq
\P(X \geq n) \leq C\rho^n$, de modo que para $\abs t < \frac1\rho$ la
serie $\sum p_n\abs t^n$ está dominada por la serie geométrica
convergente $C\sum(\rho\abs t)^n$: el radio es al menos
$\frac1\rho > 1$. El radio de la \omterm{ex:b2:powerseries:fibonacci}{función generatriz} y el decaimiento
geométrico de la cola son dos caras de una misma propiedad.
\end{solution}

\begin{solution}{exo:b2:genfun:9}
Escríbanse $p_k^{(n)} = \P(X_n = k)$ y $p_k = \P(X = k)$.

\emph{Caso $k = 0$.} Para $t \in \intoo{0}{1}$ y cualquier \omterm{def:b2:randomvar:law}{ley}
$(q_j)$ con $\sum_j q_j \leq 1$:
\[
\Bigl|\,q_0 - \sum_j q_jt^j\Bigr|
= \sum_{j \geq 1} q_j t^j
\leq \sum_{j\geq1}t^j = \frac{t}{1 - t} .
\]
De ahí,
\[
\abs{p_0^{(n)} - p_0}
\leq \frac{2t}{1 - t}
+ \abs{G_{X_n}(t) - G_X(t)} .
\]
Dado $\varepsilon > 0$, elíjase $t$ con $\frac{2t}{1-t} <
\frac\varepsilon2$ y después $n_0$ tal que el último término sea $<
\frac\varepsilon2$ para $n \geq n_0$: así, $p_0^{(n)} \to p_0$.

\emph{Paso de inducción.} Supongamos $p_j^{(n)} \to p_j$ para $j <
k$. Considérense las funciones \emph{desplazadas}
\[
g_n(t) = \frac{G_{X_n}(t) - p^{(n)}_0}{t}
= \sum_{j\geq0} p^{(n)}_{j+1}t^j,
\qquad
g(t) = \frac{G_X(t) - p_0}{t} ,
\]
\omterm{ex:b2:powerseries:fibonacci}{funciones generatrices} de las sucesiones subprobabilísticas
$(p^{(n)}_{j+1})_j$ (masa total $\leq 1$, que es todo lo que usaba el
argumento del caso $k = 0$). Para $t \in \intoo{0}{1}$ fijo, $g_n(t)
\to g(t)$ por hipótesis y por el caso $k = 0$. Aplicando el argumento
del caso $k = 0$ a $g_n$ resulta $p_1^{(n)} \to p_1$; e iterando el
desplazamiento $k$ veces se obtiene $p_k^{(n)} \to p_k$ para todo
$k$. (Esta es la instancia discreta y elemental del teorema de
\omterm{def:b2:metric:continuity}{continuidad} de Lévy, cuya forma general ---para funciones
características--- es un hito del tercer año.)
\end{solution}

\begin{solution}{exo:b2:genfun:10}
Punto a punto, $\frac{1 + (-1)^X}{2}$ vale $1$ cuando $X$ es par y
$0$ cuando es impar, de modo que, tomando \omterm{def:b2:randomvar:expectation}{esperanzas}
(transferencia),
\[
\P(X \text{ par}) = \frac{1 + \E\bigl((-1)^X\bigr)}2 =
\frac{1 + G_X(-1)}2 .
\]
Poisson: $\frac{1 + \eu^{-2\lambda}}2 \to \frac12$ cuando $\lambda$
crece. Binomial: $\frac{1 + (1 - 2p)^n}2$. En ambos casos, $G_X(-1)
\to 0$ dice que la paridad de $X$ se vuelve una moneda equilibrada:
la \omterm{def:b2:randomvar:law}{ley} se reparte sobre muchos enteros y olvida su paridad.
\end{solution}

\begin{solution}{exo:b2:genfun:11}
$t + \dots + t^6 = t\,\frac{1 - t^6}{1 - t}$ y $1 - t^6 = (1 - t)(1 +
t)(1 + t + t^2)(1 - t + t^2)$, lo que da la factorización enunciada.
Para el primer dado, $(1 + t)(1 + t + t^2) = 1 + 2t + 2t^2 + t^3$, de
modo que $\frac{t(1+t)(1+t+t^2)}6 = \frac{t + 2t^2 + 2t^3 + t^4}6$:
caras $\{1, 2, 2, 3, 3, 4\}$. Para el segundo, desarrollando
\[
(1 + 2t + 2t^2 + t^3)(1 - t + t^2)^2 = 1 + t^2 + t^3 + t^4 +
t^5 + t^7,
\]
resulta $\frac{t(1+t)(1+t+t^2)(1-t+t^2)^2}6 = \frac{t + t^3 + t^4 +
t^5 + t^6 + t^8}6$: caras $\{1, 3, 4, 5, 6, 8\}$. El producto de las
dos \omterm{ex:b2:powerseries:fibonacci}{funciones generatrices} reagrupa los seis factores en
$\bigl(\frac{t(1+t)(1+t+t^2)(1-t+t^2)}6 \bigr)^2$, el cuadrado de la
función del dado estándar: el par de Sicherman tiene exactamente la
\omterm{def:b2:randomvar:law}{ley} estándar para el total; las \omterm{ex:b2:powerseries:fibonacci}{funciones generatrices} clasifican
todas esas reagrupaciones.
\end{solution}

\begin{solution}{exo:b2:genfun:12}
Sean $A$ y $B$ las \omterm{ex:b2:powerseries:fibonacci}{funciones generatrices} de la duración restante
arrancando desde ``sin cara actual'' y desde ``con una cara actual''.
Se gasta un lanzamiento y entonces: desde el estado $0$, una cruz
devuelve al estado $0$ y una cara pasa al estado $1$; desde el estado
$1$, una cara termina el juego y una cruz devuelve al estado $0$:
\[
A(t) = t\bigl(q\,A(t) + p\,B(t)\bigr),
\qquad
B(t) = t\bigl(p + q\,A(t)\bigr).
\]
Sustituyendo: $A(1 - qt) = pt\,B = pt(pt + qtA)$, de donde
\[
G_T(t) = A(t) = \frac{p^2t^2}{1 - qt - pq\,t^2} .
\]
En $t = 1$ el denominador vale $1 - q - pq = p(1 - q) = p^2$:
$G_T(1) = 1$, el juego termina casi seguramente (como mostraba el~\cref{exo:b2:proba:6} por recursión). Derivación logarítmica en $1$:
$\E(T) = 2 - \frac{D'(1)}{D(1)}$ con $D(t) = 1 - qt - pqt^2$ y
$D'(1) = -q - 2pq$:
\[
\E(T) = 2 + \frac{q + 2pq}{p^2} = \frac{2p^2 + q + 2pq}{p^2}
= \frac{1 + p}{p^2},
\]
que vale $6$ para $p = \frac12$.
\end{solution}

\begin{solution}{pb:b2:genfun:1}
\textbf{1.} Para $t \in \intoo01$, la regla de la cadena sobre $G_n =
G \circ G_{n-1}$ da $G_n'(t) = G'\bigl(G_{n-1}(t)\bigr)G_{n-1}'(t)$.
Cuando $t \to 1^-$, $G_{n-1}(t) \uparrow 1$, y $G'$ es no decreciente
con límite por la izquierda $m$ en $1$, de modo que el primer factor
tiende a $m$; y por inducción, el segundo tiende a $m^{n-1}$. Por el~\cref{thm:b2:genfun:moments}, $\E(Z_n) = G_n'(1^-) = m^n$.

\textbf{2.} Derivando una vez más,
\[
G_n'' = G''(G_{n-1})\,(G_{n-1}')^2 +
G'(G_{n-1})\,G_{n-1}'',
\]
y haciendo $t \to 1^-$: $a_n = G''(1)m^{2(n-1)} + m\, a_{n-1}$ con
$a_n = G_n''(1)$ y $a_1 = G''(1)$. Para $m \neq 1$ se comprueba por
inducción que $a_n = G''(1)\,m^{n-1} \frac{m^n - 1}{m - 1}$ (la
recursión suma $G''(1)m^{2n-2}$ a $m\cdot
G''(1)m^{n-2}\frac{m^{n-1}-1}{m-1}$, y $m^{n-1} +
\frac{m^{n-1}-1}{m-1} = \frac{m^n - 1}{m-1}$); para $m = 1$, $a_n =
a_{n-1} + G''(1) = n\,G''(1)$.

\textbf{3.} $V(Z_n) = a_n + m^n - m^{2n}$ y $G''(1) = \sigma^2 + m^2
- m$. Para $m \neq 1$, la pieza $(m^2 -
m)m^{n-1}\frac{m^n-1}{m-1} = m^n(m^n - 1)$ cancela exactamente $m^n -
m^{2n}$, y queda $V(Z_n) = \sigma^2m^{n-1}\frac{m^n-1}{m-1}$. Para $m
= 1$: $V(Z_n) = nG''(1) = n\sigma^2$.

\textbf{4.} $Z_n$ es una variable entera no negativa, de modo que
$\P(Z_n > 0) = \P(Z_n \geq 1) \leq \E(Z_n) = m^n$ por Markov
(\cref{thm:b2:randomvar:markov}). Para $m < 1$ esto decae
geométricamente ---y de manera sumable---, así que Borel--Cantelli da
incluso que solo un número finito de generaciones son no vacías, que
es de nuevo la extinción.

\textbf{5.} Cauchy--Schwarz: $\E(Z_n)^2 = \E(Z_n\mathbf 1_{Z_n>0})^2
\leq \E(Z_n^2)\,\P(Z_n > 0)$. Con la pregunta 3 y $m < 1$:
\[
\E(Z_n^2) = V(Z_n) + m^{2n}
\leq \frac{\sigma^2m^{n-1}}{1-m} + m^{2n},
\]
de modo que, dividiendo $m^{2n}$ por esta cota y simplificando por
$m^n$,
\[
\P(Z_n > 0) \geq \frac{m^n}{\frac{\sigma^2}{m(1-m)} + m^n}
\geq \Bigl(\frac{\sigma^2}{m(1-m)} + 1\Bigr)^{-1}m^n ,
\]
usando $m^n \leq 1$ en el denominador. Con la pregunta 4: $\P(Z_n >
0) \asymp m^n$.

\textbf{6.} $G(t) = q\sum_k(pt)^k = \frac{q}{1 - pt}$, y $m = G'(1) =
\frac{pq}{(1-p)^2} = \frac pq$. Subcrítico para $p < \frac12$,
crítico para $p = \frac12$ y supercrítico para $p > \frac12$.

\textbf{7.} $G(t) = t$ se lee $pt^2 - t + q = 0$, con raíces
$\frac{1 \pm \abs{p - q}}{2p}$, es decir, $1$ y $\frac qp = \frac1m$.
La probabilidad de extinción es el punto fijo más pequeño de
$\intcc01$ (\cref{thm:b2:genfun:extinction}): $q_{\mathrm{ext}} = 1$
si $m \leq 1$, y $\frac1m$ si $m > 1$.

\textbf{8.} Para $m \neq 1$, con $p = \frac m{m+1}$ y $q =
\frac1{m+1}$: si $q_n = \frac{m^n - 1}{m^{n+1} - 1}$, entonces
\[
1 - p\,q_n = \frac{(m+1)(m^{n+1} - 1) - m(m^n - 1)}
{(m+1)(m^{n+1} - 1)} = \frac{m^{n+2} - 1}{(m+1)(m^{n+1} -
1)},
\]
luego $q_{n+1} = \frac{q}{1 - pq_n} = \frac{m^{n+1} - 1}{m^{n+2} -
1}$; y el caso base $q_0 = 0$ se cumple. Para $m = 1$: $G(t) =
\frac1{2 - t}$ y $q_{n+1} = \frac1{2 - \frac{n}{n+1}} =
\frac{n+1}{n+2}$, con $q_0 = 0$.

\textbf{9.} $1 - q_n = \frac{m^n(m - 1)}{m^{n+1} - 1}$. Para $m < 1$
el denominador tiende a $-1$: $1 - q_n \sim (1 - m)\,m^n$. Para $m >
1$:
\[
q_{\mathrm{ext}} - q_n = \frac1m - \frac{m^n - 1}{m^{n+1} -
1} = \frac{m - 1}{m\,(m^{n+1} - 1)} \sim \frac{m -
1}{m^{2}}\;m^{-n} .
\]
Y $G'(t) = \frac{pq}{(1 - pt)^2}$ evaluada en $t = \frac qp$ (donde
$1 - pt = 1 - q = p$) da $G'(q_{\mathrm{ext}}) = \frac qp =
\frac1m$: la razón observada $m^{-1}$ es exactamente la derivada en
el punto fijo atractor.

\textbf{10.} Para $p = \frac12$: $G''(t) = \frac{1/4}{(1 - t/2)^3}$,
de modo que $G''(1) = 2$ y $\sigma^2 = G''(1) + m - m^2 = 2$. La
\omterm{def:b2:multint:exact}{forma cerrada} da $1 - q_n = \frac1{n+1}$: la probabilidad de
supervivencia decae como $1/n$; demasiado despacio para ser
\omterm{def:b2:series:summable}{sumable}, a diferencia de cualquier ritmo subcrítico.

\textbf{11.} Inducción: $G_1(t) = \frac1{2-t}$ casa con la fórmula
para $n = 1$, y
\[
G(G_n(t)) = \cfrac{1}{2 - \cfrac{n - (n-1)t}{n+1 - nt}}
= \frac{n + 1 - nt}{2(n+1) - 2nt - n + (n-1)t}
= \frac{n+1 - nt}{n + 2 - (n+1)t} .
\]
Entonces
\[
\frac{G_n(t) - q_n}{1 - q_n}
= (n+1)\,\Bigl(\frac{n - (n-1)t}{n+1 - nt} -
\frac{n}{n+1}\Bigr)
= \frac{t}{n + 1 - nt}
= \frac{\frac{t}{n+1}}{1 - \frac{n}{n+1}t} ,
\]
la \omterm{ex:b2:powerseries:fibonacci}{función generatriz} de la \omterm{def:b2:randomvar:law}{ley} geométrica
$\mathcal G\bigl(\frac1{n+1} \bigr)$ sobre $\N^*$
(\cref{ex:b2:genfun:classical}): dada la supervivencia, $\P(Z_n = k
\mid Z_n > 0) = \frac1{n+1}\bigl(\frac n{n+1}\bigr)^{k-1}$, con media
condicionada $n + 1$. La media incondicionada $1 = \E(Z_n)$ es el
producto de una probabilidad de supervivencia evanescente por un
tamaño condicionado que crece linealmente.

\textbf{12.} Si el linaje se extingue en la generación $n$, entonces
$Y = Z_0 + \dots + Z_{n-1}$ es finito; y si no se extingue nunca, $Y
\geq \sum_n 1 = \infty$. Así pues, $\{Y < \infty\}$ es el \omterm{def:b2:proba:space}{suceso} de
extinción y $\P(Y < \infty) = q_{\mathrm{ext}}$. La deducción de
$H(t) = tG(H(t))$ del~\cref{exo:b2:genfun:7} ---el antepasado aporta
el factor $t$, y sus hijos fundan copias independientes de $Y$
contabilizadas por $G$--- solo usó el~\cref{thm:b2:genfun:compound}, válido en todos los regímenes.

\textbf{13.} Con $G(s) = \frac{1 + s^2}2$, la ecuación se lee $tH^2 -
2H + t = 0$, luego $H = \frac{1 - \sqrt{1 - t^2}}{t}$ (la raíz con
$H(0) = 0$). Comparando con la serie de Catalan $C(x) = \frac{1 -
\sqrt{1 - 4x}}{2x}$ (\cref{ex:b2:powerseries:catalan}): $H(t) =
\frac t2\,C\bigl(\frac{t^2}4\bigr) =
\sum_{k\geq0}C_k\,\frac{t^{2k+1}}{2^{2k+1}}$, es decir, $\P(Y = 2k+1)
= C_k2^{-2k-1}$. Comprobaciones: $\P(Y = 1) = C_0/2 = \frac12$ (el
antepasado no tiene hijos); $\P(Y = 3) = C_1/8 = \frac18$ (dos hijos,
ambos sin descendencia: $\frac12\cdot\frac12\cdot \frac12$).

\textbf{14.} Derivando $H = tG(H)$ sobre $\intoo01$ y haciendo $t \to
1^-$ (límites monótonos como en el~\cref{thm:b2:genfun:moments}): $H'(1)\bigl(1 - G'(H(1))\bigr) =
G(H(1))$. En el caso subcrítico, $H(1) = 1$ y $\E(Y) = H'(1) =
\frac1{1 - m}$. En el caso crítico, $G'(1) = 1$ anula el factor de la
izquierda mientras que el miembro derecho vale $1$: no puede existir
ningún $H'(1)$ finito, luego $\E(Y) = \infty$; y sin embargo $\P(Y <
\infty) = q = 1$.

\textbf{15.} $C_k = \frac1{k+1}\binom{2k}k \sim
\frac{4^k}{\sqrt\pi\,k^{3/2}}$ por el~\cref{ex:b2:comparison:centralbinomial}, de modo que
\[
\P(Y = 2k+1) = \frac{C_k}{2\cdot4^{k}} \sim
\frac1{2\sqrt\pi\,k^{3/2}} .
\]
Sumando la cola (por comparación con $\int_K^\infty k^{-3/2}\dd k =
2K^{-1/2}$, por arriba y por abajo): $\P(Y > 2K) \asymp K^{-1/2}$, es
decir, $\P(Y > n) \asymp n^{-1/2}$; una cola pesada con media
infinita, que cuantifica la pregunta 14.

\textbf{16.} Ambos objetos críticos ---el tiempo de retorno del paseo
equilibrado (el problema de fin de semana del~\cref{ch:b2:proba}) y
la descendencia total crítica--- son finitos casi seguramente y de
media infinita, con leyes locales de exponente $-3/2$ y colas de
exponente $-1/2$. No es casualidad: explorar un árbol genealógico
hijo a hijo produce un camino de $\pm1$ (un paso arriba por
nacimiento, uno abajo por muerte) que es exactamente un paseo
equilibrado, e $Y$ se convierte en un tiempo de primer paso. La
criticidad significa deriva nula: el proceso está siempre al borde
tanto de la extinción como de la explosión, y las fluctuaciones a
escala $\sqrt{}$ de una aleatoriedad sin deriva producen precisamente
esos exponentes.

\textbf{17.} $G''(t) = \sum_{n\geq2}n(n-1)p_nt^{n-2}$ tiene términos
no negativos, así que es no decreciente sobre $\intco01$ con límite
finito $G''(1) = \sigma^2$ (la criticidad hace $\E Z_1(Z_1 - 1) =
\sigma^2$); y una función no decreciente cuyo límite coincide con el
valor en la frontera es \omterm{def:b2:metric:continuity}{continua} en $1$. Taylor con resto integral
en el punto $1$:
\[
G(t) = 1 + (t - 1) + \int_1^t(t - s)G''(s)\,\dd s
= t + \frac{G''(1)}2(1-t)^2 + o\bigl((1-t)^2\bigr),
\]
puesto que $G''(s) = G''(1) + o(1)$ cuando $s \to 1^-$.

\textbf{18.} Reduciendo a común denominador, $h(t) = \frac{G(t) -
t}{(1 - G(t))(1 - t)}$. Por la pregunta 17, el numerador es $b(1-t)^2
+ o((1-t)^2)$ y $1 - G(t) = (1 - t)\bigl(1 - b(1-t) + o(1-t)\bigr)$,
de modo que $h(t) \to b$.

\textbf{19.} Por la definición de $h$ en $t = q_j$ y $G(q_j) =
q_{j+1}$: $\frac1{1 - q_{j+1}} - \frac1{1-q_j} = h(q_j)$; sumando
desde $j = 0$ ($q_0 = 0$) se obtiene la fórmula mostrada. Como el
proceso crítico se extingue, $q_j \uparrow 1$, luego $h(q_j) \to b$ y
la media de Cesàro $\frac1n\sum_{j<n}h(q_j) \to b$: $\frac1{1-q_n}
\sim bn$, es decir,
\[
\P(Z_n > 0) \sim \frac1{bn} = \frac{2}{\sigma^2 n} .
\]

\textbf{20.} Caso geométrico crítico: $\sigma^2 = 2$ (pregunta 10),
de modo que Kolmogórov predice $1 - q_n \sim \frac1n$; y el valor
exacto es $\frac1{n+1}$.

\textbf{21.} Como $Z_n\mathbf 1_{Z_n > 0} = Z_n$, $\E(Z_n \mid Z_n >
0) = \frac{\E(Z_n)}{\P(Z_n > 0)} = \frac1{1 - q_n} \sim
\frac{\sigma^2n}2$. En el caso geométrico esto vale $n + 1$, lo que
casa exactamente con la pregunta 11 ($\sigma^2 = 2$). La imagen
crítica: la extinción es cierta, el tamaño medio está congelado en
$1$, y los raros linajes que sobreviven tienen tamaño creciendo
linealmente; cada factor equilibrando al otro.

\textbf{22.} Para una descendencia $\mathcal P(\lambda)$, $G(t) =
\eu^{\lambda(t-1)}$ y la probabilidad de extinción es la menor raíz
de $q = \eu^{\lambda(q-1)}$. Numéricamente: $\lambda = 1.5$ da $q
\approx 0.417$ (itérese $q \mapsto \eu^{1.5(q-1)}$: $0, 0.223, 0.312,
0.356, \dots \to 0.4172$); y $\lambda = 2$ da $q \approx 0.203$. Así
pues, un caso índice desencadena un brote grande con probabilidad del
$58\,\%$ ($\lambda = 1.5$) o del $80\,\%$ ($\lambda = 2$): probable,
no seguro. La iteración desde $q_0 = 0$ converge a la raíz
\emph{menor} porque $G$ es no decreciente: por inducción, $q_n \leq
r$ para todo punto fijo $r$, y $(q_n)$ crece (es $\P(Z_n = 0)$), de
modo que su límite es un punto fijo por debajo de todos los demás.

\textbf{23.} Los $k$ antepasados fundan árboles genealógicos
\omterm{def:b2:proba:independence}{independientes}, y la extinción total es la intersección de $k$
\omterm{def:b2:proba:space}{sucesos} de extinción \omterm{def:b2:proba:independence}{independientes}: probabilidad $q^k$. Para
$\lambda = 1.5$: que la probabilidad de brote $1 - q^k$ sea $\geq
0.99$ exige $q^k \leq 0.01$, es decir, $k \geq \frac{\ln 0.01}{\ln
0.417} \approx 5.3$: seis casos iniciales hacen el brote seguro al
$99\,\%$.

\textbf{24.} \emph{$G'(q) < 1$:} $G - \mathrm{id}$ es convexa y se
anula en $q$ y en $1$, así que es $\leq 0$ sobre $\intcc q1$; si
$G'(q) = 1$, la tangente en $q$ (que la convexidad sitúa por debajo
de $G$) forzaría $G(t) \geq t$ sobre $\intcc q1$, luego $G \equiv
\mathrm{id}$ ahí, matando todos los coeficientes $p_n$ ($n \geq 2$) y
contradiciendo $m > 1$. \emph{Convergencia geométrica:} $q_n < q$
para todo $n$ (inducción, con $G$ creciente), y el teorema del valor
medio da $q - q_{n+1} = G'(c_n)(q - q_n)$ con $c_n \in \intoo{q_n}q$,
de modo que $G'(c_n) \leq G'(q) < 1$ y $q - q_n \leq q\,G'(q)^n$.
\emph{Proceso compañero:} $\widehat G(t) = G(qt)/q =
\sum_kp_kq^{k-1}t^k$ tiene coeficientes no negativos y $\widehat G(1)
= G(q)/q = 1$: es una \omterm{ex:b2:powerseries:fibonacci}{función generatriz}; y su media es $\widehat
G'(1) = G'(q) < 1$: subcrítica. Familia geométrica: $G(t) =
\frac{q}{1-pt}$, $q_{\mathrm{ext}} = \frac qp$, y
\[
\widehat G(t) = \frac pq\cdot\frac{q}{1 - p\frac qp t}
= \frac{p}{1 - qt} :
\]
la \omterm{def:b2:randomvar:law}{ley} de descendencia geométrica con $p$ y $q$ intercambiados; el
proceso supercrítico visto sobre su \omterm{def:b2:proba:space}{suceso} de extinción es el
subcrítico especular.

\textbf{25.} La tabla: $m < 1$: $q = 1$, $\P(Z_n > 0) \asymp m^n$
(preguntas 4--5), $\E(Y) = \frac1{1-m}$, y generaciones
supervivientes de media condicionada acotada. $m = 1$: $q = 1$,
$\P(Z_n > 0) \sim \frac2{\sigma^2n}$ (Kolmogórov), $\E(Y) = \infty$
con $\P(Y > n) \asymp n^{-1/2}$, y supervivientes de tamaño $\sim
\frac{\sigma^2n}2$. $m > 1$: $q < 1$ es el punto fijo más pequeño, $q
- q_n = O(G'(q)^n)$, crecimiento $\E(Z_n) = m^n$, y, condicionado a
morir, el proceso es el compañero subcrítico (pregunta 24). Las
herramientas: la composición de \omterm{ex:b2:powerseries:fibonacci}{funciones generatrices} convirtió la
recursión de poblaciones en iteración de funciones; la convexidad
fijó la geometría de los puntos fijos; Taylor en $1^-$ convirtió las
hipótesis de momentos en desarrollos locales; y el promedio de Cesàro
extrajo el $1/n$ de Kolmogórov de una suma telescópica. El volumen
del tercer año añade la martingala $Z_n/m^n$ ---cuyo límite casi
seguro refina $\E(Z_n) = m^n$ en un ritmo de crecimiento trayectoria
a trayectoria--- y el teorema de Yaglom, la \omterm{def:b2:randomvar:law}{ley} límite que hay tras
la geometría condicionada observada en la pregunta 11.
\end{solution}
